% =====================================================================
% Mục không đánh số, đặt sau Kết luận và trước tài liệu tham khảo (đề yêu cầu có trong báo cáo).
\section*{Phân công và khai báo}

\subsection*{Bảng phân công}

Bảng~\ref{tab:contrib} tóm tắt phần việc của từng thành viên, lấy từ bảng nhận việc \file{task-assignment.xlsx} và lịch sử
commit. Minh chứng chi tiết theo từng dòng việc nằm trong file bảng nhận việc và lịch sử git của repo.

\begin{table*}[t]
\centering
\caption{Phân công đóng góp.}
\label{tab:contrib}
\footnotesize
\begin{tabular}{p{3.6cm}p{13.2cm}}
\toprule
Thành viên & Phần việc \\
\midrule
Nguyễn Đình Lộc (25C11050) &
Khung repo, trích đặc trưng CLIP, module đánh giá, chia tập VG $\cap$ COCO, baseline CLIP, kiểm tra pipeline; toàn bộ cấp
độ~1 (dựng đồ thị, mô hình GATv2, huấn luyện, ablation, giải thích, 10 truy vấn) và các bước xây mô hình chính; demo. Cấp
độ~2: huấn luyện NativE trên GPU 4090, xuất embedding, kênh \texttt{kgc}, dựng lại đặc trưng MKG-W, chạy đánh giá val/test.
Bản đo giải thích v1. Điều phối, review và merge; ráp báo cáo. \\
Ngô Trương Minh Đạt (25C15003) &
Dữ liệu MKG-W: lấy Wikidata, chia tập retrieval, sinh truy vấn che tên; \texttt{EntityLinker} và hàm điểm đồ thị của bước b;
hỗ trợ tái lập NativE. \\
Trần Đắc Khoa (25C15015) &
Phân tích kết quả và bản nháp báo cáo cấp độ~2 (tái lập, bảng val/test, ablation, độ nhạy $\alpha$, trực tiếp/gián tiếp);
ví dụ giải thích trên MKG-W. \\
Nguyễn Bùi Mẫn Nhi (25C15055) &
Cấp độ~3: định nghĩa và cài đặt fidelity, validity, sparsity (bản v2, sửa các lỗi của v1), so sánh các cách giải thích
(v3: phản thực tế, GraphMask, GStarX; ba benchmark có đáp án biết trước),
bằng chứng trên từng truy vấn, cột ``bỏ phần này'' của demo; mô hình scene ranker song song (Sinkhorn trên concept graph)
và báo cáo của nó. \\
Âu Dương Khang (25C15045) &
Cấp độ~3: MMKG tiếng Việt (chọn miền, thu thập, liên kết ảnh--thực thể, duyệt triple, truy vấn tiếng Việt, data card),
chạy chuyển giao; hỗ trợ chọn encoder đa ngữ. Cấp độ~2: bộ 25 truy vấn tự viết, script đánh giá MKG-W, notebook molab;
huấn luyện NativE 3 seed và chạy lại độc lập toàn bộ cấp~2 trên molab. \\
\bottomrule
\end{tabular}
\end{table*}

\subsection*{Khai báo mã nguồn mở, dữ liệu và công cụ AI}

\textbf{Mã nguồn mở.} Nhóm dùng các thư viện qua API, không sửa mã nguồn của thư viện: PyTorch, PyTorch Geometric
(\texttt{GATv2Conv})~\cite{fey2019pyg}, open\_clip, FAISS, spaCy, SceneGraphParser, FactualSceneGraph, transformers, NLTK
(WordNet~\cite{miller1995wordnet}), Gradio. Mô hình xếp lại của cấp độ~1 do nhóm viết theo thiết kế riêng, không dùng mã
nguồn của bài báo nào. Ở cấp độ~2, nhóm chạy mã nguồn công khai của NativE (\texttt{zjukg/NativE}, commit \texttt{c72291c})
và chỉ sửa cách build: thư viện \texttt{Base.so} kèm sẵn bị lỗi segfault nên được build lại từ mã nguồn. Kênh \texttt{kgc}
tính lại hàm điểm của NativE trong code của nhóm và có test đối chiếu với bản gốc. Một lỗi của FactualSceneGraph được vá
trong code của nhóm, không sửa thư viện.

\textbf{Dữ liệu.} Visual Genome 1.4 (CC BY 4.0), caption MS-COCO 2014 (CC BY 4.0), split Karpathy, MKG-W và bộ ảnh của nhóm
tác giả MMRNS, Wikidata (CC0). Nguồn và giấy phép của MMKG tiếng Việt ghi ở Mục~\ref{sec:vn-data} và Phụ lục~\ref{app:vnmmkg}. Link tải và checksum:
\file{data/README.md}, \file{data/MANIFEST.md}.

\textbf{Công cụ AI.} Nhóm dùng Anthropic Claude (Claude Code) và OpenAI Codex để lên phương án cùng người làm, viết spec,
viết phần lớn code theo spec đã duyệt, review code, chạy thực nghiệm và soạn nháp tài liệu. Người làm của từng phần chốt
thiết kế và mọi lựa chọn (trong đó có chọn mô hình chính trước khi chấm test), duyệt code và kết quả, và chịu trách nhiệm
giải thích phần của mình. Mọi thay đổi nằm trong lịch sử git, phần lớn qua pull request có người duyệt; số liệu chỉ lấy từ
file kết quả có log và seed. Ghi chú bổ sung về cách dùng công cụ ở một số phần việc: \file{report/declaration-ai-oss.md}.
